\documentclass[10pt,a4paper]{amsart}
\usepackage[margin=19mm]{geometry}
\usepackage{amsmath,amssymb,mathtools}
\usepackage[scr=rsfs,cal=euler]{mathalfa}
\usepackage{xfrac}
\usepackage[unicode,hidelinks,bookmarks=false]{hyperref}
\newcommand{\ie}{i.e.\ }
\newcommand{\cf}{cf.\ }
\newcommand{\resp}{respectively\ }
\newcommand{\Subsection}{\subsection}
\providecommand{\nobreakdash}{\nobreak\hbox{-}}
\setlength{\emergencystretch}{2em}
\multlinegap0pt
\usepackage{amssymb}
\usepackage{mathtools}
\usepackage[shortlabels]{enumitem}
\newtheorem{lemma}{Lemma}
\newtheorem{theorem}{Theorem}
\title{Beyond Bass Collapse: New Irregular Edge-Space\\ Invariants in Ihara Theory}
\author{Hartosh Singh Bal}
\date{Source: arXiv:2604.20578v1 (CC BY 4.0)}
\begin{document}
\maketitle

\noindent\emph{Source and licence.} Beyond Bass Collapse: New Irregular Edge-Space\\ Invariants in Ihara Theory. Version arXiv:2604.20578v1 (CC BY 4.0), distributed by arXiv under the Creative Commons Attribution 4.0 International licence, http://creativecommons.org/licenses/by/4.0/, as recorded in the arXiv licence metadata for this exact version and shown on the abstract page https://arxiv.org/abs/2604.20578v1. Full licence notice: \texttt{licenses/arxiv-2604.20578-CC-BY-4.0.txt}.\par\smallskip

\noindent\textit{Editorial scope.} Counted unit: the article's theorem thm:sector (source0.tex lines 212--223). The article's complete original proof of it is delivered with it (source0.tex lines 225--328, one \texttt{proof} environment of 104 lines). No mathematics is written, repaired or re-derived: the statement and the proof are the article's own lines, in the article's own notation, and no retained line is substituted or re-flowed.  Retained uncounted as the source-local context the proof needs: lines 163--173.  Nothing was omitted from any retained span, so the package carries no omission record.  No retained line contains a citation, so the wrapper carries no bibliography and no bibliography entry.  No retained line contains a figure environment, an image-inclusion command or drawing code, so no image is delivered and the package declares no asset dependency.  Editorial formatting only, in addition: an AMS \texttt{amsart} preamble that restates the article's own declarations and macros used by the retained lines, namely the environments \texttt{lemma}, \texttt{theorem} on separate counters, together with the standard packages those lines need.  The retained environments print in this document as Lemma 1, Theorem 1 in order of appearance, whereas the article prints the same items with the numbering of its own full document.  The complete native source of the article as distributed in the arXiv e-print for this exact version is bundled unchanged as \texttt{sources/198950/source0.tex}.
\begin{lemma}[Line graph and antisymmetric line graph from incidence]\label{lem:LSfrominc}
Let $I$ denote the $m\times m$ identity. Then
\[
L := |D|^\top|D| - 2I
\]
is the adjacency matrix of the line graph $L(G)$ (in the edge-indexing of $E$), and
\[
S := D^\top D - 2I
\]
is the signed adjacency matrix of the antisymmetric line graph $\mathcal A(G)$ in the same edge-indexing. In particular, $S$ is well-defined up to switching.
\end{lemma}

\begin{theorem}[Sector decomposition]\label{thm:sector}
With $L,S$ as in Lemma~\ref{lem:LSfrominc}, one has
\[
Q^\top T Q
=
\begin{pmatrix}
\frac12 L & -\frac12\,|D|^\top D\\[4pt]
\frac12\,D^\top|D| & -\frac12 S
\end{pmatrix}.
\]
Equivalently, the restriction of $T$ to the symmetric sector is $\frac12 L$, the restriction to the antisymmetric sector is $-\frac12 S$, and the cross-coupling is the mixed incidence product.
\end{theorem}

\begin{proof}
Fix a reference orientation of each undirected edge and write the two directed realizations
of an undirected edge $e$ as $e^+$ (the reference direction) and $e^-=(e^+)^{-1}$.
Define the orthonormal $(\pm)$-basis vectors
\[
\phi_e^+ := \frac{1}{\sqrt2}\bigl(\mathbf e_{e^+}+\mathbf e_{e^-}\bigr),
\qquad
\phi_e^- := \frac{1}{\sqrt2}\bigl(\mathbf e_{e^+}-\mathbf e_{e^-}\bigr),
\]
so that $P\phi_e^\pm=\pm\phi_e^\pm$.

For $\alpha,\beta\in\{+,-\}$ define block entries
\[
T_{\alpha\beta}(e,f):=\langle \phi_e^\alpha,\,T\phi_f^\beta\rangle.
\]
Expanding $\phi_e^\alpha,\phi_f^\beta$ gives the universal identity
\begin{equation}\label{eq:block-expand}
T_{\alpha\beta}(e,f)
=
\frac12\Bigl(
T_{e^+,f^+}+\beta\,T_{e^+,f^-}+\alpha\,T_{e^-,f^+}+\alpha\beta\,T_{e^-,f^-}
\Bigr).
\end{equation}

\medskip
\noindent\textbf{Local structure: exactly one contributing directed pair.}
Assume $e\neq f$ share a vertex $w$ (otherwise all four Hashimoto entries in
\eqref{eq:block-expand} are $0$). Write $e^+=(a\to b)$ and $f^+=(c\to d)$.
Since $T_{(u,v),(x,y)}=1$ iff $v=x$ and $y\neq u$, and $e,f$ share only the single
vertex $w$, there is \emph{exactly one} choice among $(e^\pm,f^\pm)$ for which the head
of $f^\pm$ equals the tail of $e^\pm$ (hence $T=1$), and the non-backtracking
constraint is automatic because the other endpoints are distinct.

More precisely, depending on the position of $w$ with respect to the reference orientations,
exactly one of the following four situations occurs:

\[
\begin{array}{c|c|c|c|c}
\text{Case} & \text{Position of }w & \text{Nonzero entry} & D_{w,e} & D_{w,f}\\\hline
\mathrm{A} & \mathrm{head}(e^+)=\mathrm{tail}(f^+) & T_{e^+,f^+}=1 & +1 & -1\\
\mathrm{B} & \mathrm{tail}(e^+)=\mathrm{tail}(f^+) & T_{e^-,f^+}=1 & -1 & -1\\
\mathrm{C} & \mathrm{head}(e^+)=\mathrm{head}(f^+) & T_{e^+,f^-}=1 & +1 & +1\\
\mathrm{D} & \mathrm{tail}(e^+)=\mathrm{head}(f^+) & T_{e^-,f^-}=1 & -1 & +1
\end{array}
\]
(All other entries among $T_{e^\pm,f^\pm}$ vanish.)

\medskip
\noindent\textbf{Computation of the blocks.}
Using \eqref{eq:block-expand} and the fact that exactly one of the four terms is $1$:

\smallskip
\noindent\emph{(i) The symmetric block.}
For $\alpha=\beta=+$, \eqref{eq:block-expand} becomes
\[
T_{++}(e,f)=\tfrac12\bigl(T_{e^+,f^+}+T_{e^+,f^-}+T_{e^-,f^+}+T_{e^-,f^-}\bigr).
\]
Hence for adjacent distinct $e,f$ one has $T_{++}(e,f)=\tfrac12$, and for non-adjacent
edges $T_{++}(e,f)=0$. Since the diagonal is $0$, this is exactly $T_{++}=\tfrac12 L$.

\smallskip
\noindent\emph{(ii) The antisymmetric block.}
For $\alpha=\beta=-$,
\[
T_{--}(e,f)=\tfrac12\bigl(T_{e^+,f^+}-T_{e^+,f^-}-T_{e^-,f^+}+T_{e^-,f^-}\bigr).
\]
For example, in Case A the unique nonzero entry is $T_{e^+,f^+}=1$.
In \eqref{eq:block-expand} with $\alpha=\beta=-$, this entry appears with coefficient $+1$,
so $T_{--}(e,f)=\tfrac12$. Since the table gives $D_{w,e}=+1$ and $D_{w,f}=-1$ in Case A,
we have $-\tfrac12 D_{w,e}D_{w,f}=+\tfrac12$, as claimed. The remaining cases are analogous.
In Cases A--D, substituting the unique nonzero entry gives
\[
T_{--}(e,f)= -\tfrac12\,D_{w,e}D_{w,f}.
\]
Since $(D^\top D)_{e,f}=\sum_{v}D_{v,e}D_{v,f}=D_{w,e}D_{w,f}$ for adjacent $e\neq f$
(and $(D^\top D)_{e,e}=2$), it follows that
\[
T_{--} = -\tfrac12(D^\top D-2I)= -\tfrac12 S.
\]

\smallskip
\noindent\emph{(iii) The cross blocks.}
For $(\alpha,\beta)=(+,-)$,
\[
T_{+-}(e,f)=\tfrac12\bigl(T_{e^+,f^+}-T_{e^+,f^-}+T_{e^-,f^+}-T_{e^-,f^-}\bigr),
\]
and substituting the unique nonzero entry in Cases A--D yields
\[
T_{+-}(e,f)= -\tfrac12\,D_{w,f}.
\]
But for adjacent $e\neq f$,
\[
(|D|^\top D)_{e,f}=\sum_v |D|_{v,e}D_{v,f} = |D|_{w,e}D_{w,f}=D_{w,f},
\]
so $T_{+-}=-\tfrac12|D|^\top D$.

Similarly, for $(\alpha,\beta)=(-,+)$ one obtains
\[
T_{-+}(e,f)= \tfrac12\,D_{w,e}=\tfrac12(D^\top|D|)_{e,f},
\]
hence $T_{-+}=\tfrac12 D^\top|D|$.

Collecting (i)--(iii) gives the asserted block decomposition.
\end{proof}

\end{document}
